\section{Background and Problem Formulation}
\label{sec:background}

\subsection{MiniOneRec tokenizer pipeline}

Let each item $i$ be associated with a semantic embedding $\mathbf{x}_i \in \R^d$. MiniOneRec tokenizes items by training an \rqvae{} over these item embeddings and then decoding each item into a fixed-length semantic ID $\sid(i) = [z_i^{(1)}, z_i^{(2)}, z_i^{(3)}]$ \citep{kong2025minionerec,lee2022rqvae}. During \texttt{sid-train}, the model learns residual quantizers over the semantic embedding space. During \texttt{sid-generate}, the model re-encodes colliding items with Sinkhorn-based assignment enabled so that the final exported \texttt{index.json} has much lower collision than the raw train-stage codes.

This separation is important. A tokenizer should not be judged only by train-stage collision. The actual pipeline uses \texttt{sid-generate} as part of the final tokenizer construction, so a fair comparison must examine the final generated SID rather than raw codes emitted directly by the training loop.

\subsection{Why local ambiguity matters more than unrepaired global collision}

Our diagnosis of the current semantic baseline consistently points to a local failure mode. We observed many cases where the model reaches the correct semantic prefix but fails at the final leaf. This means that the tokenizer still groups together items that are semantically close yet behaviorally distinct. We refer to this phenomenon as \emph{local leaf ambiguity}.

To make this concrete, consider an item's first two SID levels $(z_i^{(1)}, z_i^{(2)})$. The set of items sharing this two-level prefix defines a same-$l_2$ bucket. If many leaves remain under the same prefix, then the final decoder still faces a difficult local choice even when the upper routing is already correct. Our current analysis therefore treats same-$l_2$ crowding and the conditional entropy of the third-level code given the prefix as tokenizer health indicators in addition to collision.

\subsection{Fair tokenizer comparison requires upstream-aligned training}

Before testing graph-aware tokenizers, we first had to resolve a provenance problem in the reproduced MiniOneRec pipeline. A direct rerun with the current local default configuration produced extremely high train-stage and generate-stage collision, which appeared inconsistent with the official low-collision Industrial index already stored in the repository. By comparing the local code path with the upstream MiniOneRec release and then rerunning with the upstream tokenizer recipe, we found that the discrepancy came primarily from hyperparameters rather than a broken implementation.

\begin{table}[t]
  \centering
  \caption{Tokenizer provenance alignment on Industrial. The upstream-style \rqvae{} regime reproduces the low-collision official tokenizer quality, whereas the local default training configuration does not.}
  \label{tab:provenance}
  \begin{tabular}{lcccc}
    \toprule
    Setting & Epochs & Batch size & Best train collision & Final index collision \\
    \midrule
    Local default rerun & 10 & 256 & 0.9989 & 0.9940 \\
    Upstream-style rerun & 10000 & 20480 & 0.0993 & 0.00434 \\
    Official committed index & -- & -- & -- & 0.00434 \\
    \bottomrule
  \end{tabular}
\end{table}

\Cref{tab:provenance} is central to the rest of the paper. It shows that fair tokenizer comparison must use the upstream-style \rqvae{} training regime. All training-time comparisons for MGR-SID in this draft therefore adopt the same tokenizer backbone, optimizer, and training schedule as the upstream MiniOneRec recipe.

\subsection{Problem statement}

These observations motivate a more specific tokenizer question than ``should collaboration be fused into tokenization?'' We instead ask:
\begin{quote}
How should graph-structured collaborative information be injected into different SID levels so that a semantic tokenizer becomes less locally ambiguous without collapsing its hierarchy?
\end{quote}

The current MGR-SID v1 instantiation answers this question conservatively. It preserves the semantic \rqvae{} backbone and adds graph-based regularization only as a structural training signal. This allows us to test the value of hierarchy-aware graph supervision before modifying the downstream SFT or RL stages.
